% Part: first-order-logic
% Chapter: axiomatic-deduction
% Section: identity

\documentclass[../../../include/open-logic-section]{subfiles}

\begin{document}

\olfileid{fol}{axd}{ide}

\olsection{\usetoken{P}{derivation} مع \usetoken{S}{identity}}

إدخال $\eq$ في ال!!{derivation}s لا يقتضي إلا زيادة مخططات بديهية.
فأما تعريف !!{derivation} وعلاقة $\Proves$ فلا تغيير فيهما، سوى أن
البديهيات المستحدثة تدخل في جملة البديهيات المأذون بها.

\begin{defn}[بديهيات !!{identity}]
\begin{align}
\ollabel{ax:id1} & \eq[{\OLId{latin-ordinary}{t}}][{\OLId{latin-ordinary}{t}}], \\
\ollabel{ax:id2} & \eq[{\OLId{latin-ordinary}{t}}_{\OLMathNumeral{1}}][{\OLId{latin-ordinary}{t}}_{\OLMathNumeral{2}}] \lif (!B({\OLId{latin-ordinary}{t}}_{\OLMathNumeral{1}}) \lif
  !B({\OLId{latin-ordinary}{t}}_{\OLMathNumeral{2}})),
\end{align}
حيث الحدود ${\OLId{latin-ordinary}{t}}$ و${\OLId{latin-ordinary}{t}}_{\OLMathNumeral{1}}$ و${\OLId{latin-ordinary}{t}}_{\OLMathNumeral{2}}$ مغلقة، أيًّا كانت.
\end{defn}

\begin{prop}\ollabel{prop:sound}
  البديهيتان \olref{ax:id1} و\olref{ax:id2} صالحتان.
\end{prop}

\begin{proof}
  يُترك إثبات الصلاحية تمرينًا.
\end{proof}

\begin{prob}
بيّن صحة \olref[fol][axd][ide]{prop:sound}.
\end{prob}

\begin{prop}\ollabel{prop:iden1}
 ${\OLId{greek-variable}{Gamma}} \Proves \eq[{\OLId{latin-ordinary}{t}}][{\OLId{latin-ordinary}{t}}]$ لكل حد مغلق ${\OLId{latin-ordinary}{t}}$ ولكل مجموعة~${\OLId{greek-variable}{Gamma}}$.
\end{prop}

\begin{prop}\ollabel{prop:iden2}
  لكل حدين مغلقين ${\OLId{latin-ordinary}{t}}_{\OLMathNumeral{1}}$ و${\OLId{latin-ordinary}{t}}_{\OLMathNumeral{2}}$، إذا كان ${\OLId{greek-variable}{Gamma}} \Proves !A({\OLId{latin-ordinary}{t}}_{\OLMathNumeral{1}})$ و
  ${\OLId{greek-variable}{Gamma}} \Proves \eq[{\OLId{latin-ordinary}{t}}_{\OLMathNumeral{1}}][{\OLId{latin-ordinary}{t}}_{\OLMathNumeral{2}}]$، فإن ${\OLId{greek-variable}{Gamma}} \Proves !A({\OLId{latin-ordinary}{t}}_{\OLMathNumeral{2}})$.
\end{prop}

\begin{proof}
ال!!{formula}
\[
(\eq[{\OLId{latin-ordinary}{t}}_{\OLMathNumeral{1}}][{\OLId{latin-ordinary}{t}}_{\OLMathNumeral{2}}] \lif (!A({\OLId{latin-ordinary}{t}}_{\OLMathNumeral{1}}) \lif !A({\OLId{latin-ordinary}{t}}_{\OLMathNumeral{2}})))
\]
من حالات~\olref{ax:id2}؛ فإذا أجرينا~\MP مرتين حصل المطلوب.
\end{proof}

\end{document}
